% Source: MQ20a_v1.html

\documentclass[11pt]{article}
\usepackage[margin=1in]{geometry}
\usepackage{amsmath}
\usepackage{amssymb}
\usepackage{siunitx}
\usepackage{enumitem}

\title{MQ20a: Kinetic Theory / Ideal Gas}
\author{}
\date{}

\begin{document}
\maketitle

\section*{MQ20a: Kinetic Theory / Ideal Gas}

\begin{enumerate}[leftmargin=*,label=\textbf{Question \arabic*.}]

\item \hfill \textit{1 pts}

A large lawnmower engine has a power output of 101 kW with a thermal efficiency of 0.57. How much heat energy (in kJ) is discarded every second?

\item \hfill \textit{1 pts}

The figure shows an engine cycle's \emph{pV}-diagram, where \emph{p}$_{1}$ = 102 kPa, \emph{p}$_{2}$ = 181 kPa, \emph{p}$_{3}$ = 297 kPa; \emph{V}$_{1}$ = 1 m$^{3}$, and \emph{V}$_{2}$ = 3.6 m$^{3}$. Calculate the net work in kJ done by the engine during one cycle.

 [Image: ME20\_PV\_5.png]

\item \hfill \textit{1 pts}

A refrigerator delivers 85 J of energy into the kitchen during a typical period of operation. If the coefficient of performance is 12, calculate the energy extracted from inside the refrigerator during the same time period.

\item \hfill \textit{1 pts}

A heat pump in heating mode takes 1,056 J from the external environment while doing 558 J of work. Calculate the coefficient of performance of the heat pump in heating mode.

\item \hfill \textit{1 pts}

Suppose a heat pump in cooling mode removes 1,021 J of energy from inside the house while ejecting 1,224 J of energy outdoors. Calculate the coefficient of performance.

\item \hfill \textit{1 pts}

A diatomic gas engine contracts isothermally from \emph{V}$_{2}$ = 2.00 m$^{3}$ and \emph{p}$_{2}$ = 230 kPa to a volume \emph{V}$_{1}$ and pressure 574 kPa. The engine then follows an isobaric process to \emph{V}$_{2}$, then an isochoric process to the original pressure, \emph{p}$_{2}$. Calculate the engine's efficiency.

[Hint: Start by drawing the engine cycle on a \emph{pV}-diagram.]

\end{enumerate}

\end{document}
